\documentclass[10pt,a4paper]{amsart}
\usepackage[margin=19mm]{geometry}
\usepackage{amsmath,amssymb,mathtools}
\usepackage[scr=rsfs,cal=euler]{mathalfa}
\usepackage{xfrac}
\usepackage[unicode,hidelinks,bookmarks=false]{hyperref}
\newcommand{\ie}{i.e.\ }
\newcommand{\cf}{cf.\ }
\newcommand{\resp}{respectively\ }
\newcommand{\Subsection}{\subsection}
\providecommand{\nobreakdash}{\nobreak\hbox{-}}
\setlength{\emergencystretch}{2em}
\multlinegap0pt
\usepackage{amssymb}
\usepackage[shortlabels]{enumitem}
\usepackage{tikz}
\usepackage{float}
\newtheorem{prop}{Proposition}
\title{Real algebraic varieties and their intersection cohomology}
\author{Chris Hone}
\date{Source: arXiv:2606.09149v1 (CC BY 4.0)}
\begin{document}
\maketitle

\noindent\emph{Source and licence.} Real algebraic varieties and their intersection cohomology. Version arXiv:2606.09149v1 (CC BY 4.0), distributed by arXiv under the Creative Commons Attribution 4.0 International licence, http://creativecommons.org/licenses/by/4.0/, as recorded in the arXiv licence metadata for this exact version and shown on the abstract page https://arxiv.org/abs/2606.09149v1. Full licence notice: \texttt{licenses/arxiv-2606.09149-CC-BY-4.0.txt}.\par\smallskip

\noindent\textit{Editorial scope.} Counted unit: the article's prop propertiesrealgeoext (source0.tex lines 432--437). The article's complete original proof of it is delivered with it (source0.tex lines 439--443, one \texttt{proof} environment of 5 lines). No mathematics is written, repaired or re-derived: the statement and the proof are the article's own lines, in the article's own notation, and no retained line is substituted or re-flowed.  No further source-local context is needed: every cross-reference of every retained line resolves inside the two delivered spans.  Nothing was omitted from any retained span, so the package carries no omission record.  No retained line contains a citation, so the wrapper carries no bibliography and no bibliography entry.  No retained line contains a figure environment, an image-inclusion command or drawing code, so no image is delivered and the package declares no asset dependency.  Editorial formatting only, in addition: an AMS \texttt{amsart} preamble that restates the article's own declarations and macros used by the retained lines, namely the environments \texttt{prop} on separate counters, together with the standard packages those lines need.  The retained environments print in this document as Proposition 1 in order of appearance, whereas the article prints the same items with the numbering of its own full document.  The complete native source of the article as distributed in the arXiv e-print for this exact version is bundled unchanged as \texttt{sources/202384/source0.tex}.
\begin{prop}\label{propertiesrealgeoext}
The real geometric extension $\mathscr{E}(Y_\mathbb{R})$ on $Y(\mathbb{R})$ is a $d$-shifted self-dual complex of sheaves; there exists an isomorphism: \[\mathbb{D}\mathscr{E}(Y_\mathbb{R})\cong \mathscr{E}(Y_\mathbb{R})[d].\]
Its cohomology therefore satisfies Poincaré duality:    
\[\mathbb{H}^i(\mathscr{E}(Y_\mathbb{R}))\cong \mathbb{H}^{d-i}_c(\mathscr{E}(Y_\mathbb{R}))^\vee.\]
In particular, if $Y(\mathbb{R})$ is compact, we have the duality \[\mathscr{E}^i(Y(\mathbb{R}),\mathbb{F}_2):=\mathbb{H}^i(\mathscr{E}(Y_\mathbb{R}))\cong \mathbb{H}^{d-i}(\mathscr{E}(Y_\mathbb{R}))^\vee=:\mathscr{E}^{d-i}(Y(\mathbb{R}),\mathbb{F}_2)^\vee\]
\end{prop}

\begin{proof}
Let $f:X\to Y$ be any resolution of singularities of $Y$, and let  $\mathbf{1}_{X(\mathbb{R})}[d]\cong \omega_{X(\mathbb{R})}$ be the (shifted) fundamental class of $X(\mathbb{R})$. As $f$ is proper, this induces the isomorphism:
\[\mathbb{D}_{Y(\mathbb{R})}f_*\mathbf{1}_{X(\mathbb{R})}\cong \mathbb{D}_{Y(\mathbb{R})} f_!\mathbf{1}_{X(\mathbb{R})}\cong f_*\mathbb{D}_{X(\mathbb{R})}\mathbf{1}_{X(\mathbb{R})}\cong f_*\omega_{X(\mathbb{R})}\cong f_*\mathbf{1}_{X(\mathbb{R})}[d].\]
As taking duals and shifts commutes with restriction to an open subset, the isomorphism class of a minimal summand $\mathscr{E}(Y_\mathbb{R})$ of $f_*\mathbf{1}_{X(\mathbb{R})}$ extending $\mathbf{1}_{Y^{sm}(\mathbb{R})}$ over the smooth locus is $d$-shifted self-dual. This self-duality of $\mathscr{E}(Y_\mathbb{R})$ then implies the Poincaré duality statement via pushforward along the terminal map to a point.
\end{proof}

\end{document}
